% REVIEW DRAFT: not an input of the current published manuscript.
% This file is not proof certification; see geometry-audit-2026-11-21.json.

\section{Adapted planes, lattices, and compatible theta models}
\label{sec:sr-theta-adapted-tower}

The rank induction requires actual lower spaces and operators, chosen
before the finite character packet. We construct these choices here.
The ground data are the specified CM extension \(E/F\), unramified
at every finite place, the standard additive character, the finite
split set \(P\), and the specified selfdual lattices. At a nonsplit
finite place the residue characteristic is odd; this is exactly the
inert-place condition of the ground datum. Put
\(\mathfrak c=\mathfrak D_{E/\mathbb Q}^{-1}\).
Since \(E/F\) is unramified,
\(\mathfrak c=\mathfrak D_{F/\mathbb Q}^{-1}\mathcal O_E\);
in particular it is fixed by conjugation. At the places over \(P\)
it is the unit ideal, by the assumed unramified additive character.

\begin{srlemma}[Isotropic planes in the actual range]
\label{lem:sr-theta-rational-planes}
For the ranks and signatures in (R6), the rational Witt indices
of the two first spaces and the target are respectively
\(p,q,p+q\). Any required first or target hyperbolic plane in the
\(A\)- or \(B\)-reduction therefore exists over \(E\).
\end{srlemma}

\begin{proof}
A nonsingular Hermitian space over a nonarchimedean quadratic
local extension has anisotropic part of dimension at most two;
over a split quadratic algebra it is hyperbolic.
Consequently a local space of dimension at least \(2j+2\)
contains \(j\) hyperbolic planes. The actual inequalities
\(a'>2(p+q)+1\) and \(a>a'\) give this finite-place condition
for \(j=p,q,p+q\) in all three spaces.
At infinity the given signatures give these Witt indices exactly.

The local-global isotropy theorem for Hermitian forms now supplies
one rational hyperbolic plane whenever the displayed index is
positive. Remove it and repeat: the finite dimension and the
two infinite signatures decrease by two and by \((1,1)\),
respectively, so the same argument supplies all the planes asserted.
For the skew-Hermitian space multiply its form by any fixed nonzero
anti-invariant element of \(E\); this converts the isotropy question
to the Hermitian one. Such a scalar does not change isotropic
subspaces. The bound at infinity also proves maximality of
the asserted rational indices.

These are precisely the field Hermitian-form classification and
local-global isotropy theorem, with no integral-lattice assertion
included in their use. One classical reference proving the
local-global theorem is Kneser, \emph{Lectures on Galois Cohomology
of Classical Groups}, \S5.5 (Landherr's theorem).
\footnote{\url{https://mathweb.tifr.res.in/sites/default/files/publications/ln/tifr47.pdf}.}
The integral splitting needed here is proved next.
\end{proof}

\begin{srlemma}[Splitting the lattice along a rational isotropic line]
\label{lem:sr-theta-integral-plane}
Let \(h\) be an \(\varepsilon\)-Hermitian form, linear in its first
variable, where \(\varepsilon=1\) or \(-1\).
Let \(D=\mathcal O_E\) in the Hermitian case and \(D=\mathfrak c\)
in the skew-Hermitian case. Suppose \(\Lambda\) is \(D\)-selfdual:
\[
 \Lambda=\{v:h(v,\Lambda)\subset D\}.
\]
Given a rational isotropic line \(Ex\), one can choose \(y\) with
\(h(x,y)=1\), \(h(y,y)=0\), and a fractional ideal \(\mathfrak b\)
such that
\[
 \begin{split}
  \Lambda\cap Ex&=\mathfrak b x,\\
  \Lambda&=\mathfrak b x\ \mathbin{\perp_{\!H}}\!
     (D\overline{\mathfrak b}^{-1}y)\ \perp\ \Lambda_1.
 \end{split}                                      \tag{M1}
 \label{eq:sr-adapted-lattice-plane}
\]
Here the first two summands form one hyperbolic plane;
\(\perp_H\) denotes their mutual hyperbolic pairing, not
orthogonality to one another. The complement \(\Lambda_1\)
is \(D\)-selfdual for the restricted form.
Moreover \(x\) may be scaled so that
\(\operatorname{ord}_w\mathfrak b=0\) at every \(w\) over \(P\).
\end{srlemma}

\begin{proof}
The intersection \(\Lambda\cap Ex\) is a rank-one projective
\(\mathcal O_E\)-module, hence \(\mathfrak b x\).
It is primitive: if a nonzero integral multiple of a vector of
\(\Lambda\) lies in \(Ex\), that vector already lies in \(Ex\).
Its quotient in \(\Lambda\) is therefore torsion free and
projective over the Dedekind domain \(\mathcal O_E\).

Dualizing this primitive summand and using \(D\)-selfduality shows
that
\[
  \Lambda\longrightarrow D\overline{\mathfrak b}^{-1},
                       \qquad z\longmapsto h(z,x)
\]
is onto. The target is projective, so a section is multiplication
by a rational vector \(y\), normalized by \(h(x,y)=1\);
the harmless factor \(\varepsilon\) in \(h(y,x)=\varepsilon\)
does not change its coefficient ideal.
Thus \(D\overline{\mathfrak b}^{-1}y\subset\Lambda\).

There is an integral \(\theta\in\mathcal O_E\) with
\(\theta+\bar\theta=1\). Indeed the trace map of a finite
unramified quadratic extension of integer rings is locally
surjective at every finite prime, including the split dyadic
primes. Its image is the unit ideal of \(\mathcal O_F\),
so it is globally surjective.
Put \(k=h(y,y)\). Then \(\bar k=\varepsilon k\).
The lattice inclusion just proved gives
\[
             k\in D^{-1}\mathfrak b\overline{\mathfrak b}.
\]
Replace \(y\) by \(y-\theta kx\).
The coefficient inclusion is preserved since
\[
 (D\overline{\mathfrak b}^{-1})
    (\theta k)\subset\mathfrak b.
\]
Bilinearity, \(h(x,x)=0\), \(h(x,y)=1\), and
\(h(y,x)=\varepsilon\) give
\[
 h(y-\theta kx,y-\theta kx)
       =k-\theta k-\varepsilon\bar\theta\bar k=0.
\]
This works with the same trace-one integer for both signs
of \(\varepsilon\).

The hyperbolic projections of \(z\in\Lambda\) are
\(h(z,y)x\) and \(\varepsilon h(z,x)y\).
The two coefficient ideals in (M1) and selfduality show that
these projections belong to \(\Lambda\). Their subtraction
gives the integral orthogonal complement. If a rational
vector in that complement pairs into \(D\) with \(\Lambda_1\),
it pairs trivially with the plane and into \(D\) with all of
\(\Lambda\); selfduality places it in \(\Lambda_1\).
This proves its selfduality.

Finally choose \(s\in E^\times\) with prescribed valuations at
the finitely many primes over \(P\), by weak approximation,
and replace \(x,y\) by \(sx,\bar s^{-1}y\).
The coefficient ideal of \(x\) becomes \(s^{-1}\mathfrak b\).
Taking \(\operatorname{ord}_w s=\operatorname{ord}_w\mathfrak b\)
proves the last assertion.
\end{proof}

Choose the planes in the large first space and in \(W\) when
\(p>0\); choose them in the small first space and in \(W\)
when \(q>0\). Apply (M1) to each. The orthogonal complements
have the required lower signatures and lattice dualities.
At \(v\in P\), the two chosen vectors are primitive in their
respective split lattice projections. Complete them by a
basis of the complement. This gives precisely the specified
two columns of an adapted split basis. For an \(A\)-step place
them in positions \(0,a-1\) in the first space and the same
positions in the first \(a\)-block of \(W\); for a \(B\)-step
use positions \(0,a'-1\) in the small first space and
\(a,d-1\) in \(W\). The lower columns are those of the
orthogonal complement. These choices are finite and are
independent of \(\alpha,\beta\).

At each infinite place choose positive and negative orthonormal
frames of the complement, and the two frames furnished by
\(x,y\) on the hyperbolic plane. The corresponding Cartan
involutions preserve the complement and carry \(x\) to a
nonzero scalar multiple of \(y\). This is the exact Cartan
compatibility required by zero evaluation, cone evaluation,
and their polynomial Fock spaces. We may first replace the
original Cartans and split bases by these adapted ones.
The lattices and rational forms have not changed.
Any two Cartans of a fixed real unitary space are conjugate
by its unitary group. The genuine oscillator action of those
conjugating elements identifies their Gaussian polynomial
spaces. First-group conjugation leaves the theta integral
unchanged by change of variables; target-group conjugation
right translates both compared functions by the same element.
Thus a comparison proved with adapted Cartans transports
back to the originally prescribed Cartans and matching maps.
All the conjugating elements are fixed by the geometric data,
before the finite packet.

\begin{srlemma}[Actual Schr\"odinger models and vacua]
\label{lem:sr-theta-fixed-models}
Every space and lower complement just constructed has a rational
Darboux Schr\"odinger model, nonzero finite lattice vacua that
are standard indicators at almost every finite place, and a
Haar measure with a measurable rational fundamental domain.
The models can be related by the partial Fourier/evaluation
maps used in the four constant-term formulas.
These choices can all be made before the finite character packet.
\end{srlemma}

\begin{proof}
The tensor space of a Hermitian first space and the skew-Hermitian
target, with the specified trace alternating form, is a
nondegenerate symplectic vector space over \(F\).
Symplectic Gram--Schmidt over \(F\) gives a rational Darboux basis.
Thus rational tensor vectors have two rational coordinate blocks,
as required by the theta summation convention.

At a finite place the tensor of the two specified lattices is
selfdual for the additive-character symplectic pairing.
This follows directly from the first lattice's
\(\mathcal O_E\)-duality and the second's
\(\mathfrak D_{E/\mathbb Q}^{-1}\)-duality.
Use the same integral trace-one \(\theta\) as above to form the
lattice Heisenberg character
\[
 \gamma\longmapsto
 \psi_v\!\left(\tfrac12 B_\theta(\gamma,\gamma)\right),
 \tag{M2}\label{eq:sr-theta-vacuum-character}
\]
where \(B_\theta\) is the tensor pairing with coefficient
\(\theta\), in the fixed convention of the Weyl operator
\(\rho(\gamma)\). This is the character in the specified
lattice-vacuum eigenrelation. Its compatibility with Weyl
multiplication follows by expanding the quadratic expression
at \(\gamma+\gamma'\), using
\(\theta+\bar\theta=1\), and using integral duality for the
remaining trace term. Explicitly, in the fixed alternating
form \(B\) the relation is
\[
 \chi_\theta(\gamma+\gamma')
   =\chi_\theta(\gamma)\chi_\theta(\gamma')
                      \psi_v(-B(\gamma,\gamma')/2).
\]
This is exactly the inverse Weyl multiplier, so it defines
a genuine character of the lift of the lattice to the
Heisenberg group. The phase is independent of the trace-one
integer chosen here. Indeed if \(\theta_1-\theta_2\) is
anti-invariant and \(k=h_{X\otimes W}(\gamma,\gamma)\),
then \(k\) is anti-invariant and belongs to \(\mathfrak c\).
Thus \((\theta_1-\theta_2)k\) is in
\(F\cap\mathfrak c=\mathfrak D_{F/\mathbb Q}^{-1}\).
Its contribution to the half trace is in the annihilator
of the standard additive character. This also relates the
same eigenline to any initially specified good-model vacuum.

There is a unique one-dimensional eigenspace for this character
in the Schr\"odinger representation. One elementary way to
see existence and dimension is to sandwich the lattice between
two sufficiently separated standard coordinate lattices.
The subspace of functions supported in the larger coordinate
box and periodic under the smaller translation lattice is
finite dimensional. The Weyl operators on it form a finite
Heisenberg representation. The selfdual lattice has maximal
commuting image; its specified character occurs exactly once,
by the finite translation/modulation Fourier basis.
The finite boxes can be enlarged to contain every vector in
question. This proves the asserted eigenline in the local
Schwartz space, including residue characteristic two.
Choose a nonzero vector on each eigenline.

For almost every finite place the Darboux basis and the tensor
lattice are standard integral, \(2\) is a unit, and the character
and the additive measure are unramified. In these coordinates
\(\mathbf1_{\mathcal O_v^N}\) has precisely (M2)'s eigenrelation.
Choose it there. At the remaining finitely many places choose
any nonzero vector on the line just constructed. This gives
the required restricted product and its actual vacua; no
simultaneous choice depending on a later character is involved.

For an upper and lower tensor space, the adapted hyperbolic
decomposition identifies the symplectic reduction of the
specified isotropic subspace with the lower tensor space.
Choose its symplectic basis by restriction and complete it to
a rational Darboux basis upstairs. Evaluation on one coordinate
block and Fourier integration on its dual block give the
claimed maps. For an arbitrary earlier Darboux basis, conjugate
by the rational symplectic change of coordinates, using its
Schr\"odinger intertwiner. The local scalar may be fixed by sending
the specified lattice vacuum to the lower one, whenever that
map is nonzero. Any resulting finite product of normalization
scalars is absorbed in the infinite intertwiner so that the
global rational evaluation map retains its specified theta
summation normalization. Only finitely many places require
this adjustment; all these scalars depend on the fixed models.
The finite vacuum calculation in
Lemma~\ref{lem:sr-cone-finite-vacuum} proves this nonvanishing
and gives the exact coefficient ideals. At infinity the
adapted positive polarization gives the ordinary Gaussian
and its polynomial derivatives; the corner evaluations and
Fourier transforms preserve these polynomial spaces as proved
in the cone and zero-evaluation calculations.

Finally a reductive unitary group is unimodular and locally
compact, with its rational subgroup discrete.
A Borel fundamental domain exists for a discrete subgroup of
a second-countable locally compact group. Choose a Haar
measure and such a domain for each first group. Quotient
integration is independent of the chosen domain. Scaling
these measures only changes the geometric constants already
kept explicit in the constant-term identities.
\end{proof}

\begin{srlemma}[Local vacuum normalization and the rational theta sum]
\label{lem:sr-theta-rational-normalization}
The finite vacuum normalizations in the preceding construction
can be imposed simultaneously with the rational theta-sum
normalization of each global intertwiner and evaluation map.
The resulting global operator and every compensation scalar
are fixed before the finite character packet.
\end{srlemma}

\begin{proof}
Start with rational Darboux bases and selfdual additive measures,
without normalizing any finite intertwiner by its vacuum.
A rational symplectic change of coordinates has a theta-preserving
Schr\"odinger intertwiner. This can be checked on the three usual
generators of the rational symplectic group:
a rational linear coordinate change, multiplication by a
rational quadratic phase, and Fourier transform in a rational
coordinate subspace. The first preserves the rational index
set and has global determinant modulus one by the product
formula. The second has phase one on rational vectors because
the global additive character is trivial on \(F\).
The third preserves the rational sum by adelic Poisson summation;
with selfdual measures the additive quotient has volume one.
Symplectic elimination expresses every rational symplectic
matrix as a product of these generators. This proves the
theta-preserving normalization of the global intertwiner.
The local factors can be taken standard at almost every place.
In particular this argument fixes the global scalar instead
of appealing just to the local Stone--von Neumann theorem.

The same argument fixes the global partial evaluation maps:
use the rational coordinate projection and Fourier integrations
specified by the adapted isotropic subspace, and transport by
these theta-preserving intertwiners. Poisson summation in the
integrated coordinate block gives their precise rational
theta-summation identity. At this stage their local factors
are fixed and standard at almost every finite place.

At a remaining finite place let \(e_v^0\) be such a factor.
It sends the upper vacuum to a nonzero scalar \(m_v\) times
the prescribed lower vacuum (or to the corresponding explicitly
computed corner-lattice indicator times that vacuum).
Nonvanishing is the finite evaluation calculation already proved:
the removed lattice subspace is totally isotropic even for
the \(E\)-valued tensor form, so the character (M2) restricts
to one on it; its compact integration has no nontrivial phase
to cancel. In coordinates the result is precisely the
nonzero ideal-coset characteristic function and Haar volume
in Lemma~\ref{lem:sr-cone-finite-vacuum}.
The lower remaining lattice character is the same (M2),
so its eigenspace is the specified lower vacuum line.

Replace \(e_v^0\) by \(m_v^{-1}e_v^0\) at these finitely many
places, and replace the combined infinite factor \(e_\infty^0\)
by \((\prod_v m_v)e_\infty^0\).
The restricted tensor product is exactly unchanged:
\[
 \left((\prod_v m_v)e_\infty^0\right)
       \otimes\bigotimes_{v<\infty}(m_v^{-1}e_v^0)
             =e_\infty^0\otimes\bigotimes_{v<\infty}e_v^0.
\]
Thus both its rational theta identity and its finite vacuum
normalizations hold at once. Multiplying an infinite
intertwiner by a nonzero scalar preserves its polynomial
Gaussian domain and all Lie/compact intertwining identities.
For a corner evaluation \(m_v\) denotes its scalar coefficient;
its corner ideal is retained, not normalized away.
Changing the initially specified good-model vacuum only changes
finitely many \(m_v\), since both choices are standard indicators
almost everywhere. By the phase-independence proved above,
it does not change the eigenline or this construction.
These are all changes of fixed lattices and models; no
Hecke-character value or conductor has been used.

For the Weil action itself, changing a scalar of a Schr\"odinger
intertwiner does not change its conjugated action.
Changing a rational polarization uses Kudla's explicit
polarization coboundary (Lemma~4.2); the global product of that
coboundary is one by the same rational quadratic-phase and
Fourier product formulas. Thus the genuine action and the
theta functional remain compatible with the fixed splitting
characters. The extra infinite scalar here belongs to the
evaluation operator, and consequently to the geometric
constant in the rank recursion; it does not alter a
character-dependent local epsilon factor.
\end{proof}

The splitting characters are also exactly those required.
For a pair of dimensions \(n,d\), Kudla's unitary splitting
uses characters \(\chi_X,\chi_W\) restricting to
\(\eta_{E/F}^{\,n},\eta_{E/F}^{\,d}\) on \(\mathbb A_F^\times\)
(Kudla's Theorem~3.1); changes of polarization are compatible
by his Lemma~4.2.
\footnote{S.~Kudla, \emph{Splitting metaplectic covers of dual
reductive pairs}, Israel J.~Math.~87 (1994), 361--401,
\url{https://doi.org/10.1007/BF02773003}.}
For the large pair these are
\((\beta,\beta/\alpha)\), and for the small pair
\((\alpha,\alpha/\beta)\).
Indeed the given restrictions are
\(\alpha|_{\mathbb A_F^\times}=\eta^{a'}\) and
\(\beta|_{\mathbb A_F^\times}=\eta^a\), while
\(a-a'\equiv a+a'\pmod2\).
Removing two dimensions preserves both restrictions.
Take these local splittings with the fixed standard additive
character and selfdual additive measures. Their restricted
tensor product is rationally trivial: the local Weil-index
factors multiply to one, and the character factors multiply
to one on \(E^\times\), by the Hecke-character condition.
It is therefore the genuine adelic Weil action used in the
rational theta series, in the rational Darboux model just chosen.
Transport by Schr\"odinger intertwiners gives that action on
any of our fixed models. Scalar ambiguity is removed by the
same local lattice normalization and global rational theta
normalization as in the earlier doubling-model comparison.
This checks the character restrictions, additive sign,
rational theta normalization, and lattice convention of the
actual model simultaneously.

\begin{srproposition}[A coherent finite family of rank reductions]
\label{prop:sr-theta-coherent-tower}
Fix the ground and geometric data with signature sum \(Y\).
There is a finite tree of actual \(A\)- and \(B\)-reductions,
with adapted planes, lower lattices, Cartans, split bases,
Darboux models, vacua, measures, and evaluation maps as above.
Every choice is made before \(\alpha,\beta\) and the finite
ramification sets. Every node satisfies the same required
rank and weight conditions.
\end{srproposition}

\begin{proof}
At an \(A\)-edge the new ranks and signatures are
\[
 (a,a';p,q)\longmapsto(a-2,a';p-1,q),
\]
so \(Y_1=Y-1\) and \(\delta_1=\delta-1\).
At a \(B\)-edge they are
\[
 (a,a';p,q)\longmapsto(a,a'-2;p,q-1),
\]
so \(Y_1=Y-1\) and \(\delta_1=\delta+1\).
In both cases \(a'_1>2Y_1+1\) and
\(\delta_1>2Y_1+1\). The parity of \(a+a'\)
is preserved, and the lower total dimension decreases by two,
so the strict weight condition remains valid.
In particular all the lower first ranks remain positive
(the terminal small rank is at least two).
Apply the preceding constructions recursively. The signature
sum decreases at each edge, so the tree is finite.
It is unnecessary to identify lower nodes reached in different
orders: the comparison induction applies to every actual
geometric node with these properties.
All operations used only this fixed finite tree and ground
data. Character restrictions at every node have just been
checked, so each later allowable packet uses these same choices.
\end{proof}
